En una zona de l'espai hi ha dues càrregues elèctriques puntuals de la mateixa magnitud però de signe contrari separades 20,0~cm.

\begin{apartat}{1}
Calculeu l'energia potencial de la distribució de càrregues.
\end{apartat}

\begin{apartat}{1}
Quin treball cal fer per a separar les càrregues des d'una distància inicial de 20,0~cm fins a una distància final de 50,0~cm?
\end{apartat}

\dades{%
  Valor absolut de cada càrrega $= 1,00\ \mathrm{\mu C}$\\
  $k = \dfrac{1}{4\pi\varepsilon_0} = 8,99\times 10^{9}\ \mathrm{N\,m^2\,C^{-2}}$}
